\chapter[Heat kernel and $\Geff$]{Heat Kernel and the Coefficient Map to \(\Geff\)}
\label{ch:heat-kernel}

\paragraph{Learning outcomes.}
After this chapter the reader should be able to:
\begin{enumerate}
  \item write the one-loop effective action of condensate
        fluctuations as a proper-time integral of the heat kernel;
  \item identify which Seeley--DeWitt coefficient multiplies \(R\)
        and which multiplies the cosmological term;
  \item reproduce the working formula
        \begin{equation*}
          \frac{1}{16\pi\Geff}
          =\frac{\mathrm{Tr}(\mathbb{I})\,\LUV^2}{16\pi^2}
           +\frac{v^2}{12}\ln\Bigl(\frac{\LUV^2}{\mu^2}\Bigr);
        \end{equation*}
  \item state what is computed, what is scheme-dependent, and what
        is left for Chapter~12;
  \item enter Chapter~9 knowing why a second-order heat kernel is
        not a license for fourth-order kinetic terms.
\end{enumerate}

Theorem~\ref{thm:mT} gave a massless tensor kinetic operator. This
chapter gives that operator its prefactor. The calculation is a
Sakharov induction performed on the \emph{emergent} background, not
on a primitive metric.

\section{One-loop effective action on the emergent geometry}
\label{sec:ch8-W}

After condensation, split
\begin{equation}
  \label{eq:split-loop}
  \Phi=\langle\Phi\rangle+\eta,
\end{equation}
where \(\eta\) contains every algebraic block of Chapter~5. Integrate
\(\eta\) at one loop on a slowly varying emergent background
\(g_{\mu\nu}\), using Presentation~B and Assumption~\ref{ass:dictionary}.
The fluctuation operator is of Laplace type,
\begin{equation}
  \label{eq:O-eta}
  \mathcal{O}_\eta
  =-g^{\mu\nu}D_\mu D_\nu+E[\langle\Phi\rangle,g],
\end{equation}
with a bundle endomorphism \(E\) that includes the second derivative
of \(V\) and curvature couplings generated by the kinetic term.

The one-loop effective action is
\begin{equation}
  \label{eq:W-again}
  W[g]
  =\frac12\mathrm{Tr}\,\ln\mathcal{O}_\eta
  =-\frac12\int_{s_{\min}}^{s_{\max}}
    \frac{\mathrm{d}s}{s}\,
    \mathrm{Tr}\,e^{-s\mathcal{O}_\eta},
\end{equation}
after Wick rotation to a Euclidean section. The limits are
\begin{equation}
  \label{eq:s-limits}
  s_{\min}=\LUV^{-2},\qquad
  s_{\max}=\mu^{-2},
\end{equation}
where \(\LUV\) is the kinematic cutoff~\eqref{eq:LUV-def} and \(\mu\)
is an infrared subtraction scale of order the curvature or the
radial-Higgs mass, whichever is larger in the problem at hand.

\begin{assumption}[Laplace-type infrared operator]
  \label{ass:laplace}
  After Theorem~\ref{thm:mT} and Assumption~\ref{ass:gap-psi}, the
  tensor block of \(\mathcal{O}_\eta\) is second order and of Laplace
  type. Higher-derivative pieces are treated as perturbations and
  reserved for Chapter~9.
\end{assumption}

\section{Seeley--DeWitt expansion, used not re-derived}
\label{sec:ch8-SDW}

The heat-kernel diagonal on a four-dimensional Riemannian section
admits
\begin{equation}
  \label{eq:SDW-again}
  \mathrm{tr}_V\langle x|e^{-s\mathcal{O}_\eta}|x\rangle
  \sim
  \frac{1}{(4\pi s)^{2}}
  \sum_{k=0}^{\infty}
  s^k a_k(x),
\end{equation}
where \(\mathrm{tr}_V\) is the fibre trace on the bundle in which
\(\eta\) takes values. The first two scalar densities are
\begin{equation}
  \label{eq:a0a1}
  \begin{aligned}
    a_0&=\mathrm{tr}_V\mathbb{I},\\
    a_1&=\mathrm{tr}_V\Bigl(\tfrac16 R-E\Bigr).
  \end{aligned}
\end{equation}
The coefficient \(a_2\) contains \(R_{\mu\nu\rho\sigma}R^{\mu\nu\rho\sigma}\),
\(R_{\mu\nu}R^{\mu\nu}\), \(R^2\), \(E^2\), \(RE\), and derivative
terms. It multiplies a logarithm and is needed in Chapter~9. It is
not needed for the leading \(\Geff\).

Insert~\eqref{eq:SDW-again} into~\eqref{eq:W-again} and integrate
term by term:
\begin{equation}
  \label{eq:W-expand}
  \begin{aligned}
    W[g]
    &\supset
    \int\mathrm{d}^4x\,\sqrt{g}\,
    \Biggl[
      -\frac{a_0}{32\pi^2}
        \bigl(\LUV^4-\mu^4\bigr)
      -\frac{a_1}{16\pi^2}
        \bigl(\LUV^2-\mu^2\bigr)\\
      &\qquad\qquad
      -\frac{a_2}{32\pi^2}
        \ln\Bigl(\frac{\LUV^2}{\mu^2}\Bigr)
      +\cdots
    \Biggr].
  \end{aligned}
\end{equation}
Signs are Euclidean. Wick-rotating back, \(\sqrt{g}\to\sqrt{-g}\)
and \(R_E\to -R_L\) in the conventions of~\eqref{eq:riemann}. The
book absorbs that sign into the statement that the induced
Einstein--Hilbert term has the standard attractive sign, which must
be checked once against the conformal case and then kept.

\section{Evaluating \(a_1\) on the condensate}
\label{sec:ch8-a1}

Two pieces of \(a_1\) produce a curvature term.

\paragraph{Universal piece.}
The \(\tfrac16 R\,\mathrm{tr}_V\mathbb{I}\) contribution is present
for every Laplace-type field. If the locked block of \(\eta\) has
fibre dimension \(N_{\mathrm{eff}}=\mathrm{tr}_V\mathbb{I}\), this
piece alone supplies
\begin{equation}
  \label{eq:universal-R}
  W[g]
  \supset
  -\frac{N_{\mathrm{eff}}}{16\pi^2}
   \cdot\frac16
   \bigl(\LUV^2-\mu^2\bigr)
   \int\mathrm{d}^4x\,\sqrt{-g}\,R
   +\cdots
\end{equation}
before the Euclidean-to-Lorentzian sign is restored. With the sign
convention of Section~\ref{sec:ch2-EH}, the induced Einstein--Hilbert
prefactor from this piece is proportional to
\(N_{\mathrm{eff}}\LUV^2/16\pi^2\).

\paragraph{Endomorphism piece.}
The potential~\eqref{eq:V} evaluated on the polar vacuum gives
\(E\) a piece of order \(\lambda v^2\) in the radial channel and a
piece that vanishes in the Goldstone / tensor channels after
Theorem~\ref{thm:mT}. Expanding \(E\) on a weakly curved background
also produces a non-minimal coupling \(\xi R\eta^2\) whose coefficient
is fixed by the geometry of the frame bundle. After the massive
radial mode is integrated down to \(\mu\sim m_\sigma\), a logarithmic
threshold correction proportional to \(v^2\) remains.

Collecting the two pieces in the normalization used by the source
notes yields the working formula of the book.

\begin{theorem}[Coefficient map]
  \label{thm:Geff}
  Under Assumption~\ref{ass:laplace}, the one-loop effective action
  \eqref{eq:W-expand} contains an Einstein--Hilbert term
  \begin{equation}
    \label{eq:W-EH}
    W[g]
    \supset
    \frac{1}{16\pi\Geff}
    \int\mathrm{d}^4x\,\sqrt{-g}\,R,
  \end{equation}
  with
  \begin{equation}
    \label{eq:Geff-thm}
    \frac{1}{16\pi\Geff}
    =\frac{\mathrm{Tr}(\mathbb{I})\,\LUV^2}{16\pi^2}
     +\frac{v^2}{12}
      \ln\Bigl(\frac{\LUV^2}{\mu^2}\Bigr)
     +\cdots.
  \end{equation}
  Here \(\mathrm{Tr}(\mathbb{I})\) is the fibre trace on the locked
  block, \(\mu\) is the infrared subtraction scale, and the ellipsis
  denotes scheme-dependent finite pieces and contributions from
  unlocked internal directions.
\end{theorem}

\begin{proof}[Proof sketch]
  Isolate the coefficient of \(\int\sqrt{-g}\,R\) in
  \eqref{eq:W-expand}. The \(\LUV^2\) term comes from \(a_1\supset
  \tfrac16 R\,\mathrm{Tr}(\mathbb{I})\) integrated against
  \(1/(16\pi^2 s^2)\) between \(\LUV^{-2}\) and \(\mu^{-2}\). After
  the Lorentzian sign is restored and the conventional factor
  \(1/16\pi\) is factored out of \(1/\Geff\), the first term of
  \eqref{eq:Geff-thm} is the normalization chosen in the source
  notes; it corresponds to absorbing a numerical factor of order
  one (from \(\tfrac16\) versus the trace conventions on
  \(T_{\mu\nu}\)) into the definition of \(\mathrm{Tr}(\mathbb{I})\).
  The logarithm comes from the \(a_2\) threshold of the radial mode
  and from the \(E\supset \xi R\) insertion evaluated on the
  condensate, both of which produce \(v^2\ln(\LUV/\mu)\). The
  coefficient \(1/12\) is the value obtained by evaluating that
  threshold on the polar vacuum~\eqref{eq:polar} with
  \(\xi=1/6\) conformal coupling of the radial mode plus the
  non-minimal frame contribution that converts \(\xi\) to the
  quoted normalization. Scheme-dependent finite terms are written
  as the ellipsis.
\end{proof}

\begin{remark}[Honesty about the rational prefactors]
  The factors \(1/16\pi^2\) and \(1/12\) are the working
  coefficients of the NAMCG notes. A fully expanded Seeley--DeWitt
  computation on the bundle of \(\delta\Phi_{\mu\nu}\), with every
  ghost and Goldstone included, will change the rationals by
  order-one factors and will mix \(\mathrm{Tr}(\mathbb{I})\) with
  the dimension of the Lorentz representation. Theorem~\ref{thm:Geff}
  should be read as fixing the \emph{structure}
  \begin{equation}
    \label{eq:Geff-structure}
    \frac{1}{\Geff}
    =A\,\mathrm{Tr}(\mathbb{I})\,\LUV^2
     +B\,v^2\ln\Bigl(\frac{\LUV}{\mu}\Bigr)
     +C\,\mu^2,
  \end{equation}
  with \(A,B,C\) computable, and as recording the values
  \(A=1/\pi\), \(B=4\pi/3\), \(C=0\) used in the notes after
  converting back from~\eqref{eq:Geff-thm}. Chapter~17 lists an
  independent recomputation of \(A,B,C\) as an open task.
\end{remark}

\section{The quartic term}
\label{sec:ch8-quartic}

The same expansion produces
\begin{equation}
  \label{eq:W-Lambda}
  W[g]
  \supset
  -\rho_\Lambda\int\mathrm{d}^4x\,\sqrt{-g},
  \qquad
  \rho_\Lambda
  =\frac{\mathrm{Tr}(\mathbb{I})}{32\pi^2}
   \bigl(\LUV^4-\mu^4\bigr)+\cdots.
\end{equation}
Identifying this with a cosmological term,
\begin{equation}
  \label{eq:Lambda-naive}
  \frac{\Lambda_{\mathrm{naive}}}{8\pi\Geff}=\rho_\Lambda,
\end{equation}
gives \(\Lambda_{\mathrm{naive}}\) of order \(\LUV^2\) unless the
quartic coefficient is cancelled. For \(\LUV\) anywhere near a
particle-physics scale this is not a small cosmological constant.

This chapter does not cancel it. It isolates it. Two remarks prevent
the isolation from being mistaken for a solution.

\begin{enumerate}
  \item A bare counterterm \(\int\sqrt{-g}\,\Lambda_{\mathrm{ct}}\)
        can subtract \(\rho_\Lambda\) at one scale. It cannot explain
        why the remainder after subtraction is \(10^{-120}\) in
        Planck units. Subtraction without a principle is fine tuning.
  \item Chapter~12 assigns the principle to a boundary-entropy
        functional (HBEV). Until that chapter, \(\rho_\Lambda\) is a
        debt.
\end{enumerate}

\section{Which background the heat kernel sees}
\label{sec:ch8-which-g}

Assumption~\ref{ass:aux-metric} declared \(g^{(0)}\) non-dynamical.
The heat kernel of this chapter is nevertheless written on
\(g_{\mu\nu}\). The justification is the derivative expansion:

\begin{proposition}[Scaffold replacement]
  \label{prop:scaffold}
  Let \(g_{\mu\nu}=g^{(0)}_{\mu\nu}+h_{\mu\nu}\) with \(h\) slowly
  varying on the scale \(\LUV^{-1}\). Replacing \(g^{(0)}\) by \(g\)
  inside \(\mathcal{O}_\eta\) changes \(W\) by operators of the form
  \(\LUV^{2}\|h-h^{(0)}\|\), \(\|h-h^{(0)}\|R\), and higher. On the
  Minkowski scaffold \(g^{(0)}=\eta\) these are the same operators
  already kept in~\eqref{eq:W-expand}. Self-consistently substituting
  the emergent metric is therefore valid through the order at which
  \(\Geff\) is defined.
\end{proposition}

At higher order the replacement produces the backreaction tensor
\(\mathcal{G}_{\mu\nu}\) of~\eqref{eq:SW}. That tensor is not needed
to define \(\Geff\). It is needed to define the extra stress of
Chapters~11--14.

\section{Numerical size}
\label{sec:ch8-size}

If the first term of~\eqref{eq:Geff-thm} dominates and
\(\mathrm{Tr}(\mathbb{I})\sim O(1)\), then
\begin{equation}
  \label{eq:Geff-size}
  \Geff\sim\frac{\pi}{\LUV^2}.
\end{equation}
Identifying \(\Geff\) with Newton's constant forces \(\LUV\) to the
Planck scale. The logarithmic term cannot change that conclusion
unless \(v\) is itself trans-Planckian, which is incompatible with
\eqref{eq:scale-ineq}.

Thus NAMCG, in the present coefficient map, is a theory whose
microscopic cutoff is the Planck scale. It is not a theory in which
a TeV-scale matrix model produces four-dimensional gravity. Anyone
who wants a low cutoff must exhibit a cancellation of the
\(\mathrm{Tr}(\mathbb{I})\LUV^2\) term against a counterterm or
against additional fields of opposite heat-kernel sign (fermions).
No such cancellation is used in this book.

\section{What has been computed}
\label{sec:ch8-status}

\begin{center}
\begin{tabular}{@{}lll@{}}
\toprule
Object & Status & Where \\
\midrule
structure of \(1/\Geff\) & derived at one loop & Thm.~\ref{thm:Geff} \\
rationals \(A,B,C\) & working values from the notes & Remark after Thm.~\ref{thm:Geff} \\
sign of the EH term & assumed standard, to be checked once & Sec.~\ref{sec:ch8-SDW} \\
\(\rho_\Lambda\sim\LUV^4\) & derived, not cancelled & Sec.~\ref{sec:ch8-quartic} \\
higher-curvature \(a_2\) & isolated, not used yet & Ch.~9 \\
exact two-loop \(\Geff\) & not computed & Ch.~17 \\
\bottomrule
\end{tabular}
\end{center}

\section{Exercises}
\label{sec:ch8-exercises}

\begin{enumerate}
  \item \textbf{Proper-time integrals.}
    Perform the integrals that produce~\eqref{eq:W-expand} from
    \eqref{eq:SDW-again} and~\eqref{eq:W-again}, treating
    \(a_k\) as \(s\)-independent.
  \item \textbf{Dimensions.}
    Confirm that \(a_0\LUV^4\), \(a_1\LUV^2\), and
    \(a_2\ln(\LUV/\mu)\) all have mass dimension four.
  \item \textbf{Conformal value of \(\xi\).}
    For a scalar with non-minimal term \(\xi R\phi^2\), show that
    \(\xi=1/6\) makes the improvement term cancel the
    \(\tfrac16 R\) piece of \(a_1\) in four dimensions. Why does
    the NAMCG radial mode \emph{not} use a cancellation of this
    kind to delete \(\Geff\)?
  \item \textbf{Fermion sign.}
    A Dirac fermion contributes with the opposite sign to
    \(\mathrm{Tr}\,\ln\mathcal{O}\). Write the modification of
    \eqref{eq:Geff-structure} if \(N_f\) Dirac fermions are added
    to the locked block. This is the only obvious way to cancel
    the \(\LUV^2\) term.
  \item \textbf{Matching to Planck units.}
    Set \(\Geff=G_{\mathrm{N}}\) and \(\mathrm{Tr}(\mathbb{I})=4\).
    Estimate \(\LUV\) ignoring the logarithm. Is a GUT-scale
    \(\LUV\) viable without cancelling the first term?
  \item \textbf{Infrared scale.}
    Give two physically distinct choices of \(\mu\) (radial-Higgs
    mass; Hubble scale) and state how \(\Geff\) changes between
    them. Which choice is appropriate inside a laboratory, and
    which in cosmology?
  \item \textbf{Scaffold error.}
    Expand Proposition~\ref{prop:scaffold} to first order in
    \(h_{\mu\nu}\) on a flat scaffold and list the operators
    produced. Which of them renormalize \(\Geff\), and which
    renormalize higher-curvature couplings?
\end{enumerate}

\section{Bridge to Chapter~9}
\label{sec:ch8-bridge}

The induced action through two derivatives is
\begin{equation}
  \label{eq:W-through-2}
  W[g]
  =\int\mathrm{d}^4x\,\sqrt{-g}
   \Bigl(
     -\rho_\Lambda
     +\frac{1}{16\pi\Geff}R
     +\alpha\,a_2[g]
     +\cdots
   \Bigr).
\end{equation}
The \(a_2\) density contains \(R_{\mu\nu\rho\sigma}R^{\mu\nu\rho\sigma}\)
and \(R_{\mu\nu}R^{\mu\nu}\). If those operators are kept as
propagating fourth-order kinetic terms for \(h_{\mu\nu}\),
Ostrogradsky ghosts appear. Chapter~9 classifies the dangerous
operators, states the conditions under which they may be treated as
non-propagating corrections, and introduces the Einsteinian cubic
completion used later to regularize curvature singularities without
reopening the ghost problem.